\chapter{Procalcitonin}

\section*{What Is This Test?}
The procalcitonin (PCT) test measures the serum concentration of procalcitonin, a 116-amino-acid peptide precursor of the calcium-regulating hormone calcitonin. In health, procalcitonin is produced by the C-cells of the thyroid gland and metabolized to calcitonin, with serum levels below approximately 0.1 ng/mL. In bacterial infection and sepsis, however, procalcitonin is massively induced in extraglandular tissues (liver, lung, kidney, adipose tissue, and immune cells) by bacterial endotoxin and proinflammatory cytokines (interleukin-1, tumor necrosis factor, and especially interleukin-6), and its secretion bypasses the normal proteolytic processing pathway. Serum procalcitonin therefore rises dramatically (often to values greater than 10 ng/mL) within 2--4 hours of a bacterial stimulus, peaks at 12--48 hours, and declines rapidly when the infection resolves. Viral infections and localized, non-bacterial inflammation produce little or no procalcitonin elevation. Because of this specificity for systemic bacterial infection, procalcitonin is used to distinguish bacterial from viral infection, to guide the decision to start antibiotics, to monitor the response to antibiotic therapy (serial measurement), and to de-escalate or stop antibiotics early, reducing unnecessary antimicrobial exposure \cite{schuetz2017procalcitonin}.

\section*{Alternative Names}
PCT; Serum Procalcitonin; Procalcitonin Level; Calcitonin Precursor; Procalcitonin-Guided Antibiotic Therapy

\section*{Principle}
Procalcitonin is measured by immunoassay on automated platforms using the time-resolved amplified cryptate emission (TRACE) technology (BRAHMS PCT) or an electrochemiluminescence immunoassay (ECLIA). The assay uses two monoclonal antibodies directed against different epitopes of the procalcitonin molecule: a capture antibody on a solid phase and a detection antibody labeled with a luminescent tag. In the sandwich format, the patient's procalcitonin bridges the two antibodies, and the measured signal is proportional to the procalcitonin concentration. The functional assay sensitivity (lowest level measurable with a coefficient of variation under 20\%) is approximately 0.06--0.1 ng/mL. Results are reported in ng/mL (micrograms per liter). The kinetic profile is important for interpretation: procalcitonin rises within 2--4 hours of a bacterial stimulus, peaks at 12--48 hours, and has a half-life of approximately 22--35 hours, so serial measurements (every 12--24 hours) are used to track the response to therapy.

\section*{History}
Procalcitonin was identified in the 1970s--1980s as the prohormone of calcitonin. Its clinical significance emerged in 1993, when Assicot and colleagues reported markedly elevated procalcitonin in children with severe bacterial sepsis but not in viral infections, establishing procalcitonin as a sepsis marker \cite{assicot1993high}. The first commercial immunoassay (BRAHMS PCT) was introduced in the mid-1990s, and automated assays followed in the 2000s. The sepsis definitions (Sepsis-3) and several large randomized controlled trials in the 2010s validated procalcitonin-guided antibiotic algorithms: the ProHOSP and ProACT trials showed that using procalcitonin thresholds to start and stop antibiotics reduces antibiotic exposure without increasing mortality \cite{huang2018procalcitonin}. The 2017 Surviving Sepsis Campaign guidelines recommended procalcitonin-based algorithms to discontinue antibiotics in patients with sepsis and documented the test's role in antibiotic stewardship \cite{rhodes2017surviving}.

\section*{Why This Test Is Ordered}
\begin{itemize}
  \item Differentiation of bacterial from viral infection: procalcitonin is markedly elevated in systemic bacterial infection and sepsis but remains low or only mildly elevated in most viral infections
  \item Guidance of antibiotic initiation in suspected lower respiratory tract infection, pneumonia, and sepsis: low procalcitonin supports withholding antibiotics
  \item Guidance of antibiotic discontinuation (antibiotic stewardship): serial procalcitonin measurement allows antibiotics to be stopped when levels fall below threshold (commonly less than 0.25--0.5 ng/mL) or decline by more than 80\% from peak
  \item Early diagnosis and severity assessment of sepsis: procalcitonin rises earlier than C-reactive protein and correlates with the severity of bacterial infection and the risk of progression to septic shock \cite{wacker2013procalcitonin}
  \item Monitoring the response to antibiotic therapy: persistent elevation or a rising level indicates treatment failure, an undrained source, or a superinfection
  \item Assessment of postoperative patients: sustained elevation beyond day 1--2 after surgery suggests postoperative bacterial infection rather than the physiologic inflammatory response
  \item Distinguishing bacterial from non-infectious causes of the systemic inflammatory response syndrome, such as pancreatitis, autoimmune flares, and trauma
  \item Evaluation of immunocompromised and critically ill patients, in whom clinical signs of infection are unreliable
  \item Supporting the diagnosis of neonatal sepsis when clinical findings are equivocal
\end{itemize}

\section*{Primary vs Secondary Test}
Procalcitonin is a secondary, adjunctive test. It is not a standalone diagnostic test for sepsis or infection; it is used together with clinical assessment, blood cultures, and other biomarkers. Its principal validated role is antibiotic stewardship: deciding whether to start antibiotics and when to stop them. It should be interpreted with the clinical context because elevations occur in non-infectious conditions (major surgery, trauma, burns, pancreatitis, prolonged cardiogenic shock, and some inflammatory states).

\section*{How This Test Is Performed}
A blood sample is collected by standard venipuncture into a serum separator tube or lithium heparin tube; procalcitonin is stable in serum and plasma for several days refrigerated. The sample is analyzed on an automated immunoassay platform. Because the clinical decisions are time-sensitive, procalcitonin is increasingly performed as a point-of-care or near-patient test. Serial measurements are typically obtained every 12--24 hours to follow the response to therapy. Blood cultures and other sepsis biomarkers (lactate, C-reactive protein, white blood cell count) are usually collected at the same time. Results are available within 1--3 hours on central analyzers and within approximately 20--30 minutes on point-of-care devices.

\section*{What Preparation Is Needed}
No fasting or special preparation is required. The result must be interpreted in the context of the timing of the suspected infection and of the sampling relative to any surgical procedure or trauma, because procalcitonin rises physiologically for the first 24--48 hours after major surgery. Medication history should be reviewed, because immunosuppressants (including corticosteroids) may blunt the response, and some drugs can cause spurious elevations or interference.

\section*{Contraindications and Precautions}
\begin{itemize}
  \item There are no absolute contraindications to standard venipuncture. Important interpretive cautions:
  \item (1) Procalcitonin is not specific for bacterial infection: elevated levels occur after major surgery, trauma, burns, severe pancreatitis, cardiogenic shock, prolonged cardiac arrest, and in some autoimmune or oncologic conditions.
  \item (2) Renal impairment does not markedly alter procalcitonin, unlike some other markers, but interpretation in acute kidney injury should be cautious.
  \item (3) Immunosuppressed patients (including those on corticosteroids or chemotherapy) may have blunted procalcitonin responses, and a normal level does not exclude serious infection.
  \item (4) A low procalcitonin does not rule out localized bacterial infection (e.g., abscess, empyema) or infections that are primarily localized rather than systemic.
  \item (5) In early infection (within the first 2--4 hours), procalcitonin may not yet be elevated; repeat sampling may be needed.
  \item (6) Procalcitonin does not reliably distinguish bacterial sepsis from certain non-bacterial causes of a massive systemic inflammatory response.
  \item (7) Heterophilic antibodies and rheumatoid factor can interfere with immunoassays, causing spurious results.
  \item (8) The optimal thresholds for starting and stopping antibiotics differ by clinical scenario and algorithm; the interpretation should follow a validated procalcitonin-guided protocol \cite{schuetz2011procalcitonin}.
\end{itemize}
\section*{Risks}
Minimal risk. Slight pain or bruising at the venipuncture site. Rarely: hematoma, infection, or vasovagal syncope.

\section*{What Happens After the Test}
The procalcitonin value is applied to a clinical algorithm. In suspected lower respiratory tract infection or sepsis, a low value (typically less than 0.25 ng/mL, or less than 0.1--0.5 ng/mL depending on the setting) supports withholding or stopping antibiotics, while a value above the threshold supports starting or continuing them. When antibiotics are started, serial procalcitonin is measured every 12--24 hours; antibiotics are discontinued when the level falls below the stopping threshold or drops by more than 80--90\% from the peak, provided the clinical course is improving. A rising or persistently elevated procalcitonin despite therapy prompts a search for an undrained source (abscess), a superinfection, or inadequate antimicrobial coverage. The test is used together with blood cultures, lactate, imaging, and clinical assessment; antibiotic decisions are never made on the procalcitonin value alone.

\section*{Understanding Results}

\subsection*{Below Normal (Less Than 0.1--0.25 ng/mL)}
\begin{itemize}
  \item \textbf{Low procalcitonin (commonly less than 0.1 ng/mL):} Systemic bacterial infection is unlikely, supporting withholding antibiotics in suspected viral or non-bacterial illness (e.g., most upper respiratory tract infections, viral pneumonia, and non-infectious inflammation).
  \item \textbf{Interpretive caution:} A low level does not exclude localized bacterial infection (abscess, empyema, cholangitis), very early infection, or infection in immunocompromised patients; clinical judgment remains paramount.
  \item \textbf{During follow-up:} A fall below the stopping threshold (typically less than 0.25--0.5 ng/mL, or greater than 80\% decline from peak) supports antibiotic discontinuation.
\end{itemize}

\subsection*{Above Normal (Greater Than 0.5 ng/mL)}
\begin{itemize}
  \item \textbf{0.5--2 ng/mL:} Suggests a bacterial infection that may be localized or evolving; the exact clinical response depends on the setting (community-acquired pneumonia, sepsis, or a non-infectious cause such as surgery or pancreatitis).
  \item \textbf{Greater than 2 ng/mL, and especially greater than 10 ng/mL:} Strongly suggests systemic bacterial infection or sepsis with a high likelihood of positive blood cultures; a very high level correlates with the severity of sepsis and the risk of septic shock.
  \item \textbf{Persistent elevation or rising levels on therapy:} Indicates treatment failure, an undrained focus, a resistant organism, or a new infection; prompts repeat cultures, imaging, and review of antimicrobial coverage.
  \item \textbf{Non-infectious causes of elevation:} Major surgery (peaks at 24--48 hours), severe trauma or burns, acute pancreatitis, prolonged cardiogenic shock, heatstroke, and some autoimmune or oncologic conditions.
  \item \textbf{Prolonged cardiogenic shock} may cause elevated procalcitonin in the absence of infection through gut wall hypoperfusion.
\end{itemize}

\section*{Normal Results}
\begin{itemize}
  \item \textbf{Healthy adults:} less than 0.1 ng/mL (0.1 microg/L)
  \item \textbf{Suggestive of systemic infection in appropriate clinical context:} greater than 0.5 ng/mL
  \item \textbf{Strongly suggestive of sepsis:} greater than 2 ng/mL
  \item \textbf{Critical values and algorithm thresholds:} vary by guideline and clinical scenario; validated procalcitonin-guided algorithms define starting thresholds (approximately 0.1--0.5 ng/mL in lower respiratory tract infection) and stopping thresholds (approximately 0.25--0.5 ng/mL or greater than 80\% peak decline)
  \item \textbf{Neonates:} reference values are higher in the first 48--72 hours of life (physiologic peak in the first day), and age-specific interpretation is required
  \item \textbf{Conversion:} 1 ng/mL procalcitonin = 1 microg/L
\end{itemize}

\section*{Follow-up Tests}
\begin{itemize}
  \item \textbf{Serial procalcitonin (every 12--24 hours):} The central follow-up test for antibiotic stewardship; trends rather than single values guide decisions
  \item \textbf{Blood cultures:} Obtained before antibiotics whenever sepsis is suspected; procalcitonin does not replace cultures
  \item \textbf{Lactate:} A sepsis biomarker and prognostic marker; elevated and persisting lactate indicates tissue hypoperfusion and correlates with mortality
  \item \textbf{C-reactive protein (CRP):} A complementary inflammatory marker; CRP rises later (24--48 hours) and is less specific for bacterial infection but remains widely used
  \item \textbf{Complete blood count with differential:} Leukocytosis or leukopenia supports the diagnosis of infection
  \item \textbf{Imaging:} Chest radiograph, CT, or ultrasound to identify the source (pneumonia, abscess, empyema, intra-abdominal infection)
  \item \textbf{Sputum, urine, and wound cultures and antimicrobial susceptibility testing:} Identify the pathogen and guide targeted therapy
  \item \textbf{Organ function monitoring (creatinine, liver enzymes, coagulation):} Sepsis is associated with multiorgan dysfunction
\end{itemize}

\section*{References}
\bibliographystyle{unsrtnat}
\bibliography{references}

\clearpage
\section*{Regulatory Status and Authorization Date}
This chapter describes a test category, not a specific commercial product. FDA, CDSCO, EU IVDR, and MHRA approval or registration dates therefore cannot be assigned without the assay manufacturer, product name, instrument, and intended use. For laboratory-developed tests, an individual agency approval date may not exist. See the Preface for the interpretation of regulatory status.

\clearpage
